\begin{abstract}
LLM data agents usually query databases with their user's full entitlement, and tool-level least privilege does not
narrow which rows and columns an allowed query tool can reach. We present task-scoped entitlements (TSE): a trusted
broker infers the tables, columns and rows a natural-language task needs, compiles that scope into an ephemeral
PostgreSQL role and security-barrier views owned by the user's role, and offers a bounded, audited widening request
when the scope is insufficient. In a preregistered evaluation on \nTasks{} BIRD and Spider tasks with Claude Haiku~5.5,
Claude Sonnet~5.5 and open-weight agents, under synthetic role policies, TSE scopes exposed \oerColReduction-fold
fewer columns and \oerCellReduction-fold fewer cells than the full entitlement (\oerCellReductionSuff-fold where the
scope sufficed). Execution accuracy changed by \dExHaikuTse{} (Haiku) and \dExSonnetTse{} (Sonnet) percentage points;
with a stronger broker model, by \hbDHaiku{} and \hbDSonnet{} (exploratory). Under injection naming an out-of-task
column, exfiltration fell from \exfSonnetCtxBzero\% to \exfSonnetCtxTse\% (Sonnet). An injection that made the agent
request the column succeeded in \widExfilAutoSonnet\% of runs under automatic approval, but in
\widExfilJudgeSonnet\% when a broker-side judge that sees only the original task approved widening (exploratory).
Scripted bypass attempts succeeded in \bypassSuccess{} of \bypassAttempts{} cases, while a SQL-parsing gate
enforcing the same scope admitted queries hidden in function arguments. TSE rests on scope-inference quality (sufficient for \suffTse\% of tasks) and correct database
enforcement; it does not protect against inference from in-scope data, covert channels, misuse of in-scope data, or
malicious users.
\end{abstract}
